% ==============================================================
\section{提案手法: $d$ 次元 Apriori-Window アルゴリズム}
\label{sec:method}
% ==============================================================

問題~\ref{prob:dense_detection}を解くための
\textbf{$d$ 次元 Apriori-Window アルゴリズム}を提案する。
以下の3つの高速化技術を組み合わせ（Algorithm~\ref{alg:md_apriori}）、
計算量を
$O\!\left(M \cdot G + N \cdot \prod_{k=1}^{d-1} W_k\right)$ に抑える
（第0軸を時間軸 $t$ とする）。
各技術は Apriori-Window~\cite{aprioriwindow}の
対応する1次元技術の直接拡張である（表~\ref{tab:correspondence}）。

\begin{table}[t]
\centering
\caption{Apriori-Window~\cite{aprioriwindow}と本研究の技術対応表。}
\label{tab:correspondence}
\small
\begin{tabular}{ll}
\toprule
Apriori-Window（1次元）& 本研究（$d$ 次元）\\
\midrule
二分探索 $O(\log n)$ 区間クエリ & $d$ 次元接頭辞和 $O(2^d)$ ウィンドウクエリ（新規）\\
$t$ 軸ストライド（$k=1$） & 空間イベントリスト + $t$ 軸ストライド（$k=1$）\\
候補区間交差による枝刈り & 候補セル集合交差による枝刈り（$k\geq2$）\\
密集区間（1次元連結成分） & 密集領域（$2d$-近傍連結成分）\\
\bottomrule
\end{tabular}
\end{table}

% ==============================================================
\subsection{$d$ 次元接頭辞和による $O(2^d)$ ウィンドウクエリ}\label{sec:prefix_sum}
% ==============================================================

出現グリッド $I_P$（サイズ $\prod N_k$）に対し、
各次元を $+1$ パディングした $d$ 次元接頭辞和 $\mathrm{PS}_P$（サイズ $\prod (N_k+1)$）を構築する：
\begin{equation}\label{eq:prefix_sum}
  \mathrm{PS}_P[\bm{b}]
  = \sum_{\bm{b}' \prec \bm{b}} I_P[\bm{b}'],
\end{equation}
ここで $\bm{b}' \prec \bm{b}$ は成分ごとに $b'_k < b_k$（全次元）を意味する。
$\mathrm{PS}_P$ は各次元を順番に累積和をとることで $O(G)$ で構築できる。
構築後、任意のウィンドウサポートは $2^d$ コーナーの包除原理により
\begin{equation}\label{eq:query}
  S_P(\bm{s})
  = \sum_{\bm{c} \in \{0,1\}^d}
    (-1)^{d - |\bm{c}|_1}
    \mathrm{PS}_P[\bm{s} + \bm{c} \odot \bm{W}]
\end{equation}
で $O(2^d)$ に取得できる（$\bm{c} \odot \bm{W}$ は成分積）。
$k \geq 2$ ではこれを全グリッド点に適用し、
サポートサーフェス $S_P$ を $O(2^d \cdot G)$ で計算する。
実用的な $d$（$\leq 5$ 程度）では $2^d$ は定数とみなせる。

% ==============================================================
\subsection{$t$ 軸 1 次元ストライド最適化（$k = 1$）}\label{sec:t_stride}
% ==============================================================

$k = 1$（単体アイテム）の密集セル検出では、
Apriori-Window の $t$ 軸ストライド~\cite{aprioriwindow}を
各空間ウィンドウ $(s_1, \ldots, s_{d-1})$ に独立に適用する。

\paragraph{空間イベントリストの構築.}
トランザクション $(\bm{c}_n, d_n)$（ビン化後: $\bm{b}_n$）がアイテム $i$ を含む場合、
$i$ が関係する空間ウィンドウ開始点の第 $k$ 成分（$k \geq 1$）の範囲は
\[
  s_k \in [\max(0,\, b_n^{(k)} - W_k + 1),\; b_n^{(k)}]
\]
であるから、全空間次元の組み合わせ $(s_1, \ldots, s_{d-1})$ に
$b_n^{(0)}$（時間ビン）を追加する。
構築コストは $O\!\left(N_i \cdot \prod_{k=1}^{d-1} W_k\right)$。
$d = 1$（時間のみ）のとき単一リスト $L_i() = \mathrm{occ}(\{i\})$ に帰着する。

\paragraph{$t$ 軸ストライド.}
各 $(s_1, \ldots, s_{d-1})$ について、ソート済みリスト $L_i(s_1, \ldots, s_{d-1})$ に対し
ウィンドウ左端 $l$ を動かして密集窓の開始点を求める
（Algorithm~\ref{alg:t_stride}）。
本アルゴリズムは Apriori-Window の $t$ 軸ストライドと同一の
ジャンプ規則を使用する。

ウィンドウ $[l, l+W_t)$ 内のカウント $c$ が $\theta$ 未満の場合、
次のイベントを含む最小の窓開始点へジャンプする：
\begin{equation}\label{eq:jump_empty}
  l \leftarrow \max(l+1,\; L_i(\bm{s}_{\mathrm{sp}})[\mathrm{next}] - W_t + 1).
\end{equation}
カウント $c > \theta$ の場合（余剰 $\mathrm{surplus} = c - \theta$ 個）、
左端から $\mathrm{surplus}$ 番目のイベントの時刻へジャンプする：
\begin{equation}\label{eq:jump_dense}
  l \leftarrow L_i(\bm{s}_{\mathrm{sp}})[l_{\mathrm{lo}} + \mathrm{surplus}].
\end{equation}
$l_{\mathrm{lo}}$ はリスト中で $l$ 以上の最初のインデックスである。

\begin{lemma}[ストライドの正確性]\label{lem:stride}
式~\eqref{eq:jump_dense} のジャンプ先 $l'$ に対し、
全ての $l'' \in (l, l')$ についても $c(l'') \geq \theta$ が成立する。
すなわち、スキップされた位置も全て密集窓の開始点である。
\end{lemma}

\begin{proof}
$l' = L_i[l_{\mathrm{lo}} + \mathrm{surplus}]$ とおく。
任意の $l'' \in (l, l')$ に対し、
ウィンドウ左端を $l$ から $l''$ に進めると
$L_i[l_{\mathrm{lo}}], \ldots, L_i[l_{\mathrm{lo}} + \mathrm{surplus} - 1]$
のうち $l''$ 未満の要素がウィンドウから脱出するが、
$L_i[l_{\mathrm{lo}} + \mathrm{surplus}] = l' > l''$ より
$L_i[l_{\mathrm{lo}} + \mathrm{surplus}]$ はまだウィンドウ内に残る。
したがって $c(l'') \geq c(l) - \mathrm{surplus} = \theta$。
\end{proof}

補題~\ref{lem:stride}は Apriori-Window の命題~3~\cite{aprioriwindow}の
$d$ 次元版であり、スキップ区間内の密集窓の開始点を
陽に評価せずとも補填により正確に列挙できることを示す。

\begin{algorithm}[t]
\caption{$t$ 軸 1 次元ストライドによる密集窓の開始点列挙}
\label{alg:t_stride}
\begin{algorithmic}[1]
\Require ソート済みイベントリスト $L$、ウィンドウ幅 $W_t$、閾値 $\theta$
\Ensure 密集窓の開始点の集合 $\mathcal{S}$
\State $\mathcal{S} \leftarrow \emptyset$,\quad
       $l \leftarrow L[0] - (W_t - 1)$,\quad
       $\mathit{prev} \leftarrow \mathrm{nil}$
\While{$l \leq L[\mathrm{last}]$}
  \State $l_{\mathrm{lo}} \leftarrow \mathrm{lower\_bound}(L, l)$,\quad
         $l_{\mathrm{hi}} \leftarrow \mathrm{upper\_bound}(L, l + W_t - 1)$
  \State $c \leftarrow l_{\mathrm{hi}} - l_{\mathrm{lo}}$
  \If{$c \geq \theta$}
    \If{$\mathit{prev} \neq \mathrm{nil}$}
      \State $\mathcal{S} \leftarrow \mathcal{S} \cup \{\mathit{prev}+1, \ldots, l-1\}$
             \Comment{ストライドで飛ばした区間を補填}
    \EndIf
    \State $\mathcal{S} \leftarrow \mathcal{S} \cup \{l\}$,\quad $\mathit{prev} \leftarrow l$
    \If{$c > \theta$}
      \State $l \leftarrow \max(l+1, L[l_{\mathrm{lo}} + (c - \theta)])$
             \Comment{式~\eqref{eq:jump_dense}}
    \Else
      \State $l \leftarrow l + 1$
    \EndIf
  \Else
    \State $\mathit{prev} \leftarrow \mathrm{nil}$
    \State $l \leftarrow \max(l+1, L[l_{\mathrm{hi}}] - W_t + 1)$
           \Comment{式~\eqref{eq:jump_empty}；$L[l_{\mathrm{hi}}]$ が次イベント}
  \EndIf
\EndWhile
\Return $\mathcal{S}$
\end{algorithmic}
\end{algorithm}

% ==============================================================
\subsection{Apriori 枝刈りによる候補セル限定（$k \geq 2$）}\label{sec:alg_k2}
% ==============================================================

$k \geq 2$ のパターン $P = \{i_1, \ldots, i_k\}$ について、
補題~\ref{lem:apriori} より
\begin{equation}\label{eq:candidate_cells}
  \mathcal{L}_\theta(S_P)
  \subseteq \bigcap_{Q \subset P,\, |Q|=k-1} \mathcal{L}_\theta(S_Q).
\end{equation}
右辺の交差集合を\emph{候補セル集合} $\mathcal{C}(P)$ と定義する。
これは Apriori-Window の「候補区間交差」~\cite{aprioriwindow}の $d$ 次元版であり、
密集セル集合の交差として定式化される。
$\mathcal{C}(P) = \emptyset$ ならば $P$ の密集セルは存在しないと確定し、
$I_P$ の構築・接頭辞和の計算を省略できる。

$\mathcal{C}(P) \neq \emptyset$ の場合は、
$\mathcal{C}(P)$ 内のトランザクションのみを走査して $I_P$ を構築し（$O(N)$）、
接頭辞和 $\mathrm{PS}_P$（$O(G)$）・サポート曲面（$O(G)$）を計算する。

% ==============================================================
\subsection{アルゴリズム全体と計算量}\label{sec:complexity}
% ==============================================================

\begin{algorithm}[t]
\caption{$d$ 次元 Apriori-Window アルゴリズム}
\label{alg:md_apriori}
\begin{algorithmic}[1]
\Require $d$ 次元トランザクション $\mathcal{D}$、パラメータ $(\bm{W}, \theta, L_{\max})$
\Ensure 各パターン $P$ の密集領域集合 $\{\mathcal{R}_P\}$
\State \textbf{離散化}: $\mathcal{D}$ を $(\bm{b}_n, d_n)$ に変換、グリッドサイズ $(N_0, \ldots, N_{d-1})$ を決定
\ForAll{$i \in \mathcal{I}$ \textbf{in parallel}}  \Comment{$k=1$; Rayon 並列}
  \State $L_i(\bm{s}_{1:}) \leftarrow $ 空間イベントリスト構築（\S\ref{sec:t_stride}）
  \ForAll{$\bm{s}_{1:} \in \prod_{k=1}^{d-1} [0, N_k)$}
    \If{$L_i(\bm{s}_{1:}) \neq \emptyset$}
      \State $\mathcal{S} \leftarrow $ \Call{TStrideStarts}{$L_i(\bm{s}_{1:}), W_0, \theta$}
             \Comment{Algorithm~\ref{alg:t_stride}}
      \State 各 $s_0 \in \mathcal{S}$ を密集セル集合 $D_i$ に追加
    \EndIf
  \EndFor
  \State $\mathcal{R}_{\{i\}} \leftarrow$ $2d$-近傍 Union-Find 連結成分 of $D_i$
\EndFor
\State $\mathcal{F}_1 \leftarrow \{\{i\} : \mathcal{R}_{\{i\}} \neq \emptyset\}$
\While{$\mathcal{F}_{k-1} \neq \emptyset$ \textbf{and} $k \leq L_{\max}$}  \Comment{$k \geq 2$}
  \State \texttt{cands} $\leftarrow$ \Call{GenerateCandidates}{$\mathcal{F}_{k-1}$}
  \State \texttt{cands} $\leftarrow$ \Call{PruneCandidates}{\texttt{cands}, $\mathcal{F}_{k-1}$}
  \ForAll{$P \in$ \texttt{cands} \textbf{in parallel}}
    \State $\mathcal{C}(P) \leftarrow \bigcap_{Q \subset P,|Q|=k-1} \mathcal{L}_\theta(S_Q)$
           \Comment{式~\eqref{eq:candidate_cells}}
    \If{$\mathcal{C}(P) = \emptyset$} \textbf{continue} \EndIf
    \State $I_P \leftarrow$ \Call{BuildItemGrid}{$\mathcal{D}, P, \mathcal{C}(P)$}
    \State $\mathrm{PS}_P \leftarrow$ \Call{BuildPrefixSumND}{$I_P$}
    \State $S_P \leftarrow$ \Call{ComputeSupportSurface}{$\mathrm{PS}_P, \bm{W}$}
    \State $\mathcal{R}_P \leftarrow$ $2d$-近傍 Union-Find 連結成分 of
           $\{\bm{s} : S_P(\bm{s}) \geq \theta\}$
  \EndFor
  \State $\mathcal{F}_k \leftarrow \{P : \mathcal{R}_P \neq \emptyset\}$,\quad
         $k \leftarrow k + 1$
\EndWhile
\Return $\{R_P : P \in \bigcup_k \mathcal{F}_k\}$
\end{algorithmic}
\end{algorithm}

\begin{table}[t]
\centering
\caption{$d$ 次元 Apriori-Window 各ステップの計算量。
$N$: トランザクション数、$G = \prod_{k=0}^{d-1} N_k$、
$W_{\mathrm{sp}} = \prod_{k=1}^{d-1} W_k$（空間ウィンドウ積）、$C_k$: $k$ 段候補数。}
\label{tab:complexity}
\small
\begin{tabular}{lll}
\toprule
ステップ & 計算量 & 備考 \\
\midrule
離散化 & $O(N)$ & \\
$k=1$: 空間イベントリスト構築 & $O(N \cdot W_{\mathrm{sp}})$ & 全アイテム合計 \\
$k=1$: $t$ 軸ストライド & $O(|\mathcal{S}_i| + \log N_i)$ per 空間キー & $|\mathcal{S}_i|$: 密集点数 \\
$k \geq 2$: 接頭辞和構築 & $O(2^d \cdot G)$ per 候補 & Apriori 枝刈り後 \\
$k \geq 2$: $I_P$ 構築 & $O(N)$ per 候補 & \\
連結成分 & $O(|\mathcal{L}| \cdot \alpha)$ per パターン & $\alpha$: 逆アッカーマン \\
\bottomrule
\end{tabular}
\end{table}

支配的コストは $k=1$ の空間イベントリスト構築 $O(N W_{\mathrm{sp}})$ と、
$k \geq 2$ の接頭辞和構築 $O(C_k \cdot 2^d \cdot G)$（$C_k$: Apriori 枝刈り後の候補数）である。
ナイーブ実装の $O\!\left(M \cdot G \cdot \prod_{k=0}^{d-1} W_k\right)$ に対して、
$t$ 軸ストライドはウィンドウ幅方向の削減、
Apriori 枝刈りはパターン次元の削減を担う。
次元数 $d$ が増えると $2^d$ と $G$ の双方が増加するが、
実用的な範囲（$d \leq 5$）では管理可能である（EX-E; \ref{sec:ex_e}節）。
